\begin{theorem}[三类换位碰撞边界下四个 Langlands 分量的环面秩与导子指数]\label{thm:xi-terminal-zm-s4-langlands-factor-torus-rank-conductor-table}
\leanverified{paper\_xi\_terminal\_zm\_s4\_langlands\_factor\_torus\_rank\_conductor\_table}
沿用定理 \ref{thm:xi-s4-isotypic-jacobian-toric-rank-by-boundary-type} 的三类换位碰撞边界记号 \(r\in\{1,2,3\}\)，以及结论 \ref{con:xi-terminal-zm-s4-jacobian-group-algebra-prym-e3-b3} 的 \(\QQ\)-同源分解
\[
\Jac(\widetilde C)\ \sim_{\QQ}\ J_{\mathrm{LY}}\times E_{\mathrm{res}}^2\times E^3\times B^3
\]
与等型对应
\(
\varepsilon\leftrightarrow J_{\mathrm{LY}},\
\rho_2\leftrightarrow E_{\mathrm{res}}^2,\
\rho_3\leftrightarrow E^3,\
\rho_3'\leftrightarrow B^3
\)。
对每个分量 \(A\in\{J_{\mathrm{LY}},E_{\mathrm{res}},E,B\}\)，记其在对应半稳定退化处的环面秩为 \(t(A)\)。
则对三类碰撞型 \(r\) 有表：
\[
\begin{array}{c|cccc}
r & t(J_{\mathrm{LY}}) & t(E_{\mathrm{res}}) & t(E) & t(B)\\
\hline
1 & 1 & 1 & 1 & 2\\
2 & 0 & 0 & 1 & 1\\
3 & 1 & 0 & 0 & 1
\end{array}
\]
并且在特征 \(0\) 的半稳定几何退化处（或更一般地在温和半稳定情形下），\(\ell\)-进 Tate 模的惯性表示无 Swan 部分，因而 Artin 导子指数等于环面秩。
\end{theorem}
\begin{proof}
由定理 \ref{thm:xi-s4-isotypic-jacobian-toric-rank-by-boundary-type}，在碰撞型 \(r\) 处四个 \(\QQ[S_4]\)-等型因子 \(A_\varepsilon,A_{\rho_2},A_{\rho_3},A_{\rho_3'}\) 的环面秩
\((t_\varepsilon,t_{\rho_2},t_{\rho_3},t_{\rho_3'})\)
由该定理中的表格确定。
另一方面，结论 \ref{con:xi-terminal-zm-s4-jacobian-group-algebra-prym-e3-b3} 的等型对应表明
\[
A_\varepsilon\sim_{\QQ} J_{\mathrm{LY}},\qquad
A_{\rho_2}\sim_{\QQ}E_{\mathrm{res}}^2,\qquad
A_{\rho_3}\sim_{\QQ}E^3,\qquad
A_{\rho_3'}\sim_{\QQ}B^3.
\]
环面秩对同源不变且对直积可加，故
\[
t(J_{\mathrm{LY}})=t_\varepsilon,\qquad
2\,t(E_{\mathrm{res}})=t_{\rho_2},\qquad
3\,t(E)=t_{\rho_3},\qquad
3\,t(B)=t_{\rho_3'}.
\]
将三类 \(r\) 的数值代入即得表格。
\par
最后，在特征 \(0\) 的半稳定几何退化处（或温和半稳定情形下），惯性表示的 Swan 部分为 \(0\)，而 Artin 导子指数等于环面秩与 Swan 部分之和，从而导子指数等于环面秩。证毕。
\end{proof}

\begin{corollary}[碰撞选择律与 \texorpdfstring{$B$}{B}-通道的普遍非良约化]\label{cor:xi-terminal-zm-s4-langlands-factor-collision-selection}
\leanverified{paper\_xi\_terminal\_zm\_s4\_langlands\_factor\_collision\_selection}
在定理 \ref{thm:xi-terminal-zm-s4-langlands-factor-torus-rank-conductor-table} 的设定下：
\begin{enumerate}
  \item 对所有三类换位碰撞边界均有 \(t(B)\ge 1\)，因此 \(B\) 在这些边界处均为半稳定但非良约化；并且在 \(r=1\) 型边界上有 \(t(B)=2\)。
  \item \(J_{\mathrm{LY}}\) 的良约化仅发生在 \(r=2\) 型边界；\(E_{\mathrm{res}}\) 的良约化发生在 \(r\in\{2,3\}\)；\(E\) 的良约化仅发生在 \(r=3\) 型边界。
\end{enumerate}
\end{corollary}
\begin{proof}
良约化等价于环面秩为 \(0\)。直接读取定理 \ref{thm:xi-terminal-zm-s4-langlands-factor-torus-rank-conductor-table} 的表格即可。证毕。
\end{proof}

\begin{theorem}[两类 Prym 的环面秩分层]\label{thm:xi-terminal-zm-s4-prym-torus-rank-layering}
\leanverified{paper\_xi\_terminal\_zm\_s4\_prym\_torus\_rank\_layering}
沿用结论 \ref{con:xi-terminal-zm-s4-involution-prym-polarization-eigensplitting} 的记号，
令 \(\pi_\tau:\widetilde C\to C_\tau\) 为换位商诱导的二次覆盖，\(\pi_\delta:\widetilde C\to C_\delta\) 为双换位商诱导的二次覆盖。
则在三类换位碰撞边界 \(r\in\{1,2,3\}\) 上，其 Prym 的环面秩分别为
\[
t\!\bigl(\Prym(\pi_\delta)\bigr)=(6,4,2),
\qquad
t\!\bigl(\Prym(\pi_\tau)\bigr)=(7,3,3),
\]
其中括号中三元组按 \(r=(1,2,3)\) 顺序排列。
\end{theorem}
\begin{proof}
由结论 \ref{con:xi-terminal-zm-s4-involution-prym-polarization-eigensplitting} 的 \(\QQ\)-同源分解，
\[
\Prym(\pi_\tau)\sim_{\QQ} J_{\mathrm{LY}}\times E_{\mathrm{res}}\times E\times B^2,
\qquad
\Prym(\pi_\delta)\sim_{\QQ} E^2\times B^2.
\]
环面秩对直积可加并对同源不变，故
\[
t\!\bigl(\Prym(\pi_\delta)\bigr)=2t(E)+2t(B),
\qquad
t\!\bigl(\Prym(\pi_\tau)\bigr)=t(J_{\mathrm{LY}})+t(E_{\mathrm{res}})+t(E)+2t(B).
\]
将定理 \ref{thm:xi-terminal-zm-s4-langlands-factor-torus-rank-conductor-table} 的三类 \(r\) 下 \((t(J_{\mathrm{LY}}),t(E_{\mathrm{res}}),t(E),t(B))\) 逐项代入即可得到所示数值。证毕。
\end{proof}

\endinput
